% !TEX program = xelatex
\documentclass[11pt,a4paper]{article}

\usepackage[UTF8]{ctex}
\usepackage[margin=2.2cm]{geometry}
\usepackage{amsmath, amssymb, amsthm}
\usepackage{listings}
\usepackage{xcolor}
\usepackage{tikz}
\usepackage{booktabs}
\usepackage{multirow}
\usepackage{longtable}
\usepackage{hyperref}
\usetikzlibrary{positioning, arrows.meta, shapes.geometric, fit, calc, decorations.pathreplacing}

\hypersetup{
  colorlinks=true,
  linkcolor=blue!60!black,
  urlcolor=teal!70!black,
}

\lstdefinelanguage{Rust}{
  morekeywords={fn,let,mut,pub,struct,impl,use,mod,self,Self,Option,Some,None,return,if,else,for,while,loop,match,break,continue,as,where,trait,type,enum,move,ref,const,static,unsafe,extern,crate,super,dyn,async,await,box,in},
  sensitive=true,
  morecomment=[l]{//},
  morecomment=[s]{/*}{*/},
  morestring=[b]",
}

\lstset{
  language=Rust,
  basicstyle=\ttfamily\footnotesize,
  keywordstyle=\color{blue!60!black}\bfseries,
  commentstyle=\color{green!40!black}\itshape,
  stringstyle=\color{red!60!black},
  numbers=left,
  numberstyle=\tiny\color{gray},
  numbersep=8pt,
  frame=single,
  framerule=0.4pt,
  rulecolor=\color{gray!40},
  breaklines=true,
  breakatwhitespace=true,
  showstringspaces=false,
  tabsize=4,
  columns=fullflexible,
  keepspaces=true,
}

\title{\textbf{halfedge\,: OBJ / PLY / STL / OFF / glTF 读写与网格构建模块设计文档} \\ \large \texttt{src/io.rs}}
\author{halfedge 库}
\date{2026-07-05}

\begin{document}
\maketitle
\tableofcontents

\section{模块定位}

\texttt{io.rs} 提供\textbf{五类网格文件格式}的读写能力：

\begin{center}
\renewcommand{\arraystretch}{1.18}
\begin{tabular}{@{}lll@{}}
  \toprule
  \textbf{类别} & \textbf{API} & \textbf{语义} \\
  \midrule
  \multirow{2}{*}{构建}
    & \texttt{build\_mesh\_from\_vertices\_and\_faces} & 顶点+三角面索引 $\to$ 完整半边网格 \\
    & \texttt{build\_mesh\_from\_polygons}             & 顶点+任意多边形面 $\to$ 三角化半边网格 \\
  \midrule
  \multirow{2}{*}{OBJ 读}
    & \texttt{load\_obj(path)}    & 从文件加载 \\
    & \texttt{parse\_obj(text)}   & 从文本解析 \\
  \midrule
  \multirow{2}{*}{OBJ 写}
    & \texttt{save\_obj(mesh, path)} & 保存到文件 \\
    & \texttt{format\_obj(mesh)}     & 序列化为文本 \\
  \midrule
  \multirow{4}{*}{PLY 读}
    & \texttt{load\_ply(path)}                & 从文件加载（自动判别 ASCII / 二进制） \\
    & \texttt{parse\_ply(text)}               & 从文本解析 PLY ASCII（保留旧入口） \\
    & \texttt{parse\_ply\_bytes(bytes)}       & 从字节流解析（自动判别 ASCII / 二进制） \\
    & （内部）\texttt{parse\_ply\_binary\_le} & 二进制小端 PLY 解析 \\
  \midrule
  \multirow{4}{*}{PLY 写}
    & \texttt{save\_ply(mesh, path)}            & 保存为 PLY 文件（ASCII） \\
    & \texttt{format\_ply(mesh)}                & 序列化为 PLY ASCII 文本 \\
    & \texttt{save\_ply\_binary(mesh, path)}    & 保存为 PLY 文件（二进制小端） \\
    & \texttt{format\_ply\_binary(mesh)}        & 序列化为 PLY 二进制字节流 \\
  \midrule
  \multirow{5}{*}{STL 读}
    & \texttt{load\_stl(path)}                  & 从文件加载（自动判别 ASCII / 二进制） \\
    & \texttt{parse\_stl\_bytes(bytes)}         & 从字节切片解析（启发式判别） \\
    & \texttt{parse\_stl\_ascii(text)}          & 从文本解析 STL ASCII \\
    & \texttt{parse\_stl\_binary(bytes, n)}     & 从字节解析 STL 二进制 \\
  \midrule
  \multirow{4}{*}{STL 写}
    & \texttt{save\_stl\_ascii(mesh, path)}     & 保存为 STL ASCII 文件 \\
    & \texttt{format\_stl\_ascii(mesh)}         & 序列化为 STL ASCII 文本 \\
    & \texttt{save\_stl\_binary(mesh, path)}    & 保存为 STL 二进制文件 \\
    & \texttt{format\_stl\_binary(mesh)}        & 序列化为 STL 二进制字节 \\
  \midrule
  \multirow{2}{*}{OFF 读}
    & \texttt{load\_off(path)}    & 从文件加载 OFF（ASCII） \\
    & \texttt{parse\_off(text)}   & 从文本解析 OFF \\
  \midrule
  \multirow{2}{*}{OFF 写}
    & \texttt{save\_off(mesh, path)} & 保存到 OFF 文件 \\
    & \texttt{format\_off(mesh)}     & 序列化为 OFF 文本 \\
  \midrule
  \multirow{2}{*}{glTF 读}
    & \texttt{load\_glb(path)}    & 从文件加载 GLB（glTF 二进制容器） \\
    & \texttt{parse\_glb(bytes)}  & 从字节流解析 GLB \\
  \midrule
  \multirow{2}{*}{glTF 写}
    & \texttt{save\_glb(mesh, path)} & 保存到 GLB 文件 \\
    & \texttt{format\_glb(mesh)}     & 序列化为 GLB 字节流 \\
  \midrule
  \multirow{3}{*}{统一入口}
    & \texttt{MeshError}                            & 统一错误枚举（\texttt{Obj} / \texttt{Ply} / \texttt{Stl} / \texttt{Off} / \texttt{Gltf} / \texttt{UnsupportedFormat}） \\
    & \texttt{load\_mesh(path)}                     & 按扩展名自动选择 OBJ / PLY / STL / OFF / GLB 解析器 \\
    & \texttt{save\_mesh(mesh, path)}               & 按扩展名自动选择 OBJ / PLY / STL / OFF / GLB 序列化器 \\
  \bottomrule
\end{tabular}
\end{center}

\section{支持的格式总览}

\begin{center}
\renewcommand{\arraystretch}{1.18}
\begin{tabular}{@{}llll@{}}
  \toprule
  \textbf{格式} & \textbf{扩展名} & \textbf{变体} & \textbf{特点} \\
  \midrule
  OBJ  & \texttt{.obj}            & ASCII                 & 支持 n-gon、\texttt{v/vt/vn}、负索引 \\
  PLY  & \texttt{.ply}            & ASCII + 二进制小端    & 自动判别；二进制支持任意属性类型 \\
  STL  & \texttt{.stl}            & ASCII + 二进制        & 自动判别；顶点按位模式去重 \\
  OFF  & \texttt{.off}            & ASCII                 & 简单文本；支持任意多边形面 \\
  glTF & \texttt{.glb} / \texttt{.gltf} & GLB 二进制      & 最小子集：单 primitive + POSITION + indices \\
  \bottomrule
\end{tabular}
\end{center}

\section{OBJ 格式约定}

\begin{lstlisting}
# 顶点（1-based 索引）
v 0.0 0.0 0.0
v 1.0 0.0 0.0
v 0.0 1.0 0.0
# 三角面 / n-gon（支持 v / v/vt / v/vt/vn 形式，仅取 v）
f 1 2 3
f 1/1/1 2/2/1 3/3/1
f 1 2 3 4   # 四边形 n-gon 自动 fan 三角化
# 负索引：从末尾倒数（OBJ 标准）
f -3 -2 -1
\end{lstlisting}

\begin{itemize}
  \item 仅解析 \texttt{v} 与 \texttt{f} 行；\texttt{vt}、\texttt{vn}、\texttt{g}、
        \texttt{usemtl} 等其他行被忽略；
  \item 顶点索引 1-based；负数表示从末尾倒数；
  \item 支持 n-gon：面顶点数 $\ge 3$ 即可，自动按 fan 方式三角化；
        仅当面顶点数 $< 3$ 时返回 \texttt{NotTriangular} 错误；
  \item 索引越界返回 \texttt{IndexOutOfRange} 错误。
\end{itemize}

\section{错误类型}

\texttt{ObjError} 共 4 个变体：

\begin{center}
\renewcommand{\arraystretch}{1.18}
\begin{tabular}{@{}ll@{}}
  \toprule
  \textbf{变体} & \textbf{触发场景} \\
  \midrule
  \texttt{Io(std::io::Error)}     & 文件读写错误 \\
  \texttt{Parse\{line, msg\}}     & 解析失败（行号 + 描述） \\
  \texttt{IndexOutOfRange\{line, idx, vertex\_count\}} & 面索引越界 \\
  \texttt{NotTriangular\{line, face\_verts\}} & 面顶点数 $< 3$ \\
  \bottomrule
\end{tabular}
\end{center}

\subsection{错误信息语言约定}

全模块错误信息\textbf{统一使用中文}，与代码库其余部分（\texttt{panic!}、\texttt{eprintln!}、
\texttt{expect()}、\texttt{validate.rs}）保持一致。具体包括：
\begin{itemize}
  \item 各 \texttt{Display} 实现的标题文字，如 \texttt{"IO 错误: \{e\}"}、
        \texttt{"第 \{line\} 行 PLY 解析错误: \{msg\}"}、\texttt{"OBJ 错误: \{e\}"}；
  \item \texttt{Parse\{msg\}} 中的描述串，如 \texttt{"无法解析顶点坐标"}、
        \texttt{"无效顶点坐标: \{s\}"}、\texttt{"无效 x 坐标: \{...\}"}、
        \texttt{"面索引 \{v\} 越界（顶点计数 \{v\_count\}）"}；
  \item \texttt{Unsupported} 中的描述串，如 \texttt{"不支持大端 PLY"}、
        \texttt{"未知 PLY 格式: \{other\}"}、\texttt{"不支持 GLB 版本 \{version\}（仅支持 2）"}。
\end{itemize}
\textbf{技术术语保留原文}：\texttt{OBJ} / \texttt{PLY} / \texttt{STL} / \texttt{OFF} / \texttt{glTF}
格式名、\texttt{chunk} / \texttt{bufferView} / \texttt{accessor} / \texttt{componentType} /
\texttt{POSITION} 等 GLTF 规范字段、\texttt{h.next} / \texttt{twin.prev} 等半边字段路径
不在翻译范围内，以保证错误信息与规范文档的可对照性。

\subsection{数据丢弃警告}

\begin{itemize}
  \item \texttt{build\_mesh\_from\_polygons} 跳过顶点数 $< 3$ 的退化面时，
        会通过 \texttt{eprintln!} 输出警告统计；
  \item PLY/OFF 解析器不再自行过滤退化面（移除 \texttt{if indices.len() >= 3}
        / \texttt{if k < 3} 的静默跳过），而是将所有面交给
        \texttt{build\_mesh\_from\_polygons} 统一处理与警告；
  \item 所有导出函数（\texttt{format\_obj} / \texttt{format\_ply} /
        \texttt{format\_ply\_binary} / \texttt{format\_stl\_ascii} /
        \texttt{format\_stl\_binary} / \texttt{format\_off} / \texttt{format\_glb}）
        跳过退化面/非三角面时，会通过 \texttt{eprintln!} 输出警告统计。
\end{itemize}

\section{build\_mesh\_from\_vertices\_and\_faces 算法}

\subsection{整体流程}

\begin{center}
\resizebox{\textwidth}{!}{%
\begin{tikzpicture}[
  node distance=0.5cm,
  proc/.style={draw, rounded corners=2pt, minimum width=11cm, minimum height=0.7cm, align=center, font=\small},
  op/.style={-Stealth, thick},
  startend/.style={draw, rounded corners=10pt, minimum width=2.6cm, minimum height=0.7cm, font=\small, fill=gray!12}
]
  \node[startend] (s) {开始};
  \node[proc, below=of s, fill=blue!8] (p1) {1. 创建所有顶点，记录 \texttt{v\_ids[i] $\to$ VertexId}};
  \node[proc, below=of p1, fill=blue!8] (p2) {2. 为每个面创建 3 条内部半边（CCW），设 next/prev/face，记录 \texttt{edge\_map[(i,j)] = h}};
  \node[proc, below=of p2, fill=blue!8] (p3) {3. 遍历所有有向边，查反向边：命中 $\to$ 互设 twin；未命中 $\to$ 创建边界 twin};
  \node[proc, below=of p3, fill=yellow!15] (p4) {4. 用 walking 算法设置边界半边的 next/prev};
  \node[startend, below=of p4] (e) {返回 \texttt{MeshStorage}};
  \draw[op] (s) -- (p1);
  \draw[op] (p1) -- (p2);
  \draw[op] (p2) -- (p3);
  \draw[op] (p3) -- (p4);
  \draw[op] (p4) -- (e);
\end{tikzpicture}
}
\end{center}

\paragraph{容量预分配}
入口处调用 \texttt{mesh.reserve(vertices.len(), faces.len() * 6, faces.len())}
预分配内存，其中半边容量上限 \texttt{faces.len() * 6} 来自每三角面 3 条内部半边
\texttt{+} 3 条潜在边界 twin。\texttt{build\_mesh\_from\_polygons} 同样在入口处
调用 \texttt{reserve}，半边容量按 $\sum_i k_i \times 2$ 计算（$k_i$ 为第 $i$ 个面
的边数）。预分配可将批量构建的 rehash 次数降为 0。

\subsection{步骤 1--2：顶点与内部半边}

对每个面 $[i_0, i_1, i_2]$（0-based），创建三条半边：
\[
  h_0: v_{i_0} \to v_{i_1}, \quad
  h_1: v_{i_1} \to v_{i_2}, \quad
  h_2: v_{i_2} \to v_{i_0}.
\]
设置 $\mathrm{next}(h_0)=h_1,\ \mathrm{next}(h_1)=h_2,\ \mathrm{next}(h_2)=h_0$，
$\mathrm{prev}$ 反向，$\mathrm{face} = f$。
记录 \texttt{edge\_map}：$(i_0, i_1) \to h_0$，$(i_1, i_2) \to h_1$，$(i_2, i_0) \to h_2$。

\paragraph{面索引越界检查}
\texttt{build\_mesh\_from\_vertices\_and\_faces} 与 \texttt{build\_mesh\_from\_polygons}
作为公开构建函数，在访问 \texttt{v\_ids[i]} 之前显式检查 $i < |\texttt{v\_ids}|$。
若越界，\texttt{panic} 并输出索引值与顶点总数（而非让标准库抛出匿名的
``index out of bounds''）。这把 silent panic 转为可诊断的失败，
为后续接入 \texttt{Result} 变体（缺陷 \#17）预留接口。

\subsection{步骤 3：twin 建立}

对每条有向边 $(a, b) \to h$：
\begin{itemize}
  \item 查 \texttt{edge\_map} 是否有 $(b, a)$：
    \begin{itemize}
      \item 命中 $h'$：内部边，互设 $h.\mathrm{twin} = h'$，$h'.\mathrm{twin} = h$；
      \item 未命中：边界边，创建反向边界半边 $h^*$（$\mathrm{face} = \texttt{None}$），
            互设 twin，加入 \texttt{boundary\_twins} 列表。
    \end{itemize}
\end{itemize}

\subsection{步骤 4：边界半边 next/prev 链接（walking 算法）}

\textbf{问题}：边界半边 $bh$（$A \to B$，$\mathrm{face} = \texttt{None}$）需要
设置 $\mathrm{next}$ 与 $\mathrm{prev}$，使其与其它边界半边构成外环。

\textbf{旧公式（仅适用单三角形）}：
\[
  bh.\mathrm{next} = bh.\mathrm{twin}.\mathrm{prev}.\mathrm{twin}, \quad
  bh.\mathrm{prev} = bh.\mathrm{twin}.\mathrm{next}.\mathrm{twin}.
\]
当 $bh.\mathrm{twin}.\mathrm{prev}$ 的 twin 是内部半边时，此公式失效。

\textbf{新算法（walking）}：

\textbf{求 $bh.\mathrm{next}$}：绕 $bh.\mathrm{tip}$（即 $B$）沿 CCW 方向旋转，
直到找到边界 outgoing 半边。

\begin{center}
\begin{tikzpicture}[
  node distance=0.5cm,
  proc/.style={draw, rounded corners=2pt, minimum width=10cm, minimum height=0.7cm, align=center, font=\small},
  dec/.style={diamond, draw, aspect=2.4, fill=orange!12, inner sep=1pt, font=\small, align=center, minimum width=3.5cm},
  op/.style={-Stealth, thick},
  startend/.style={draw, rounded corners=10pt, minimum width=2.6cm, minimum height=0.7cm, font=\small, fill=gray!12}
]
  \node[startend] (s) {$bh$};
  \node[proc, below=of s, fill=blue!8] (p1) {$cur \gets bh.\mathrm{twin}$（指向 $B$ 的内部半边）};
  \node[dec, below=of p1] (d1) {$p \gets cur.\mathrm{prev}$；$pt \gets p.\mathrm{twin}$；$pt.\mathrm{face} = \texttt{None}$？};
  \node[proc, right=1.6cm of d1, fill=green!12] (yes) {$bh.\mathrm{next} \gets pt$};
  \node[proc, below=of d1, fill=yellow!15] (no) {$cur \gets pt$；回到 $d1$};
  \draw[op] (s) -- (p1);
  \draw[op] (p1) -- (d1);
  \draw[op] (d1) -- node[above, font=\scriptsize] {是} (yes);
  \draw[op] (d1) -- node[right, font=\scriptsize] {否} (no);
  \draw[op, dashed] (no.east) -- ++(0.6,0) |- (d1.east);
\end{tikzpicture}
\end{center}

\textbf{求 $bh.\mathrm{prev}$}：绕 $bh.\mathrm{origin}$（即 $A$）沿 CW 方向旋转，
直到找到 twin 为边界半边的 outgoing。

\begin{center}
\begin{tikzpicture}[
  node distance=0.5cm,
  proc/.style={draw, rounded corners=2pt, minimum width=10cm, minimum height=0.7cm, align=center, font=\small},
  dec/.style={diamond, draw, aspect=2.4, fill=orange!12, inner sep=1pt, font=\small, align=center, minimum width=3.5cm},
  op/.style={-Stealth, thick},
  startend/.style={draw, rounded corners=10pt, minimum width=2.6cm, minimum height=0.7cm, font=\small, fill=gray!12}
]
  \node[startend] (s) {$bh$};
  \node[proc, below=of s, fill=blue!8] (p1) {$cur \gets bh$（$A$ 的 outgoing，twin 指向 $B$）};
  \node[dec, below=of p1] (d1) {$t \gets cur.\mathrm{twin}$；$tn \gets t.\mathrm{next}$；$tnt \gets tn.\mathrm{twin}$；$tnt.\mathrm{face} = \texttt{None}$？};
  \node[proc, right=1.6cm of d1, fill=green!12] (yes) {$bh.\mathrm{prev} \gets tnt$};
  \node[proc, below=of d1, fill=yellow!15] (no) {$cur \gets tn$；回到 $d1$};
  \draw[op] (s) -- (p1);
  \draw[op] (p1) -- (d1);
  \draw[op] (d1) -- node[above, font=\scriptsize] {是} (yes);
  \draw[op] (d1) -- node[right, font=\scriptsize] {否} (no);
  \draw[op, dashed] (no.east) -- ++(0.6,0) |- (d1.east);
\end{tikzpicture}
\end{center}

\textbf{终止性}：边界半边的 tip 必为边界顶点，边界顶点的 outgoing 环是开链
（非闭合），walking 必终止。代码设 \texttt{max\_iter = halfedge\_count + 1}
作安全阀。

\section{面绕向一致性约束}

\texttt{build\_mesh\_from\_vertices\_and\_faces} 假设所有面绕向一致
（每条无向边在两个面中方向相反）。若输入违反此约束（如两个面都包含
有向边 $a \to b$），则：

\begin{enumerate}
  \item \texttt{edge\_map} 第二次插入会覆盖第一次；
  \item 第一次插入的半边 $h_{\text{first}}$ 不在 \texttt{directed\_edges} 中，
        永远不会被设置 twin；
  \item 第二次插入的半边 $h_{\text{second}}$ 查找 $(b, a)$ 失败，创建边界 twin；
  \item 结果：$h_{\text{first}}.\mathrm{twin} = \texttt{None}$，
        \texttt{validate\_topology} 会报告 \texttt{HalfEdgeDanglingTwin}
        或入口不一致。
\end{enumerate}

\textbf{使用建议}：调用方应保证面绕向一致（如所有面 CCW 朝外）；
若不确定，应先用 \texttt{validate\_topology} 检查构建结果。

\section{OBJ 序列化}

\texttt{format\_obj} 输出格式：
\begin{itemize}
  \item 顶点按 \texttt{vertex\_ids()} 顺序，分配 1-based 索引；
  \item 面按 \texttt{face\_ids()} 顺序，每行 \texttt{f i j k}；
  \item 非三角面（边界环长度 $\ne 3$）跳过；
  \item 顶点位置用 \texttt{{:.6}} 格式化（6 位小数）。
\end{itemize}

\section{STL 读写}

STL（STereoLithography）格式广泛用于 3D 打印与 CAD 交换。
\texttt{io.rs} 支持 ASCII 与二进制两种变体，并提供自动判别。

\subsection{ASCII 格式}

\begin{lstlisting}
solid model
  facet normal 0.0 0.0 1.0
    outer loop
      vertex 0.0 0.0 0.0
      vertex 1.0 0.0 0.0
      vertex 0.0 1.0 0.0
    endloop
  endfacet
endsolid model
\end{lstlisting}

\texttt{parse\_stl\_ascii} 解析流程：
\begin{enumerate}
  \item 按行拆分，跳过 \texttt{solid} / \texttt{endsolid} / \texttt{facet normal} /
        \texttt{outer loop} / \texttt{endloop} / \texttt{endfacet} 等关键字；
  \item 提取所有 \texttt{vertex x y z} 行为顶点位置；
  \item 每连续 3 个顶点构成一个三角面；
  \item \textbf{顶点去重}：用 \texttt{f64::to\_bits()} 的位模式作为 \texttt{HashMap}
        键，相同位置的顶点复用同一 \texttt{VertexId}；
  \item 调用 \texttt{build\_mesh\_from\_vertices\_and\_faces} 构建半边网格。
\end{enumerate}

\subsection{二进制格式}

\begin{lstlisting}
[80 字节文件头]
[4 字节 little-endian 三角面数 N]
[50 字节 × N：每个三角面 12 字节法向 + 36 字节三顶点 + 2 字节属性]
\end{lstlisting}

\texttt{parse\_stl\_binary} 解析流程：
\begin{enumerate}
  \item 跳过 80 字节文件头；
  \item 读取 4 字节 \texttt{u32} 作为面数 $N$；
  \item 校验文件长度 $= 84 + 50N$，不匹配返回 \texttt{BadBinarySize} 错误；
  \item 每个三角面读取 50 字节，提取三顶点位置（法向被忽略，由网格重建）；
  \item \textbf{顶点去重}：用 \texttt{f32::to\_bits()} 的位模式（\texttt{u32} 数组）
        作为 \texttt{HashMap} 键，相同位置的顶点复用同一 \texttt{VertexId}；
  \item 调用 \texttt{build\_mesh\_from\_vertices\_and\_faces} 构建半边网格。
\end{enumerate}

\subsection{ASCII / 二进制自动判别}

\texttt{parse\_stl\_bytes(bytes)} 使用\textbf{启发式}判别：
\begin{itemize}
  \item 检查文件长度：若 $\text{len} = 84 + 50N$（$N$ 为前 80 字节后的 4 字节
        \texttt{u32}），判定为二进制；
  \item 否则尝试以 UTF-8 解码后调用 \texttt{parse\_stl\_ascii}；
  \item 两者都失败返回 \texttt{StlError::Parse}。
\end{itemize}

\textbf{为何用文件长度判别？} ASCII STL 起始 80 字节通常含 \texttt{solid} 关键字，
但二进制文件头也可能是任意 80 字节。文件长度恰好等于 $84 + 50N$ 是二进制的强信号。

\subsection{STL 序列化}

\texttt{format\_stl\_ascii} 输出标准 ASCII 格式：
\begin{itemize}
  \item \texttt{solid halfedge} 起始；
  \item 每面输出 \texttt{facet normal}（用 \texttt{face\_normal} 计算）+
        三顶点 \texttt{vertex}；
  \item \texttt{endsolid halfedge} 结束。
\end{itemize}

\texttt{format\_stl\_binary} 输出标准二进制格式：
\begin{itemize}
  \item 80 字节文件头（前缀 \texttt{halfedge binary stl}，余 0）；
  \item 4 字节 little-endian 面数 $N$；
  \item 每面 50 字节：法向（12B）+ 三顶点（36B）+ 属性（2B，全 0）。
\end{itemize}

\subsection{StlError}

\begin{center}
\renewcommand{\arraystretch}{1.18}
\begin{tabular}{@{}ll@{}}
  \toprule
  \textbf{变体} & \textbf{触发场景} \\
  \midrule
  \texttt{Io(std::io::Error)}                       & 文件读写错误 \\
  \texttt{Parse\{line, msg\}}                       & ASCII 解析失败（行号 + 描述） \\
  \texttt{BadBinarySize\{actual, expected\}}        & 二进制文件长度不匹配 \\
  \texttt{NotTriangular\{face\_verts\}}             & 非三角面（顶点数 $\ne 3$） \\
  \bottomrule
\end{tabular}
\end{center}

\subsection{顶点去重原理}

STL 文件每个三角面独立存储三顶点，无共享顶点索引。
直接构建网格会产生 $3F$ 个 \texttt{VertexId}（$F$ 为面数），
而非实际的 $V$ 个顶点（$V \approx F/2$）。

\textbf{去重算法}：
\[
  \text{key}(p) = [\texttt{f64::to\_bits}(p_x),\ \texttt{f64::to\_bits}(p_y),\ \texttt{f64::to\_bits}(p_z)]
\]
\begin{itemize}
  \item 用 \texttt{HashMap<[u64; 3], u32>}（ASCII）或
        \texttt{HashMap<[u32; 3], u32>}（二进制，因 \texttt{f32}）；
  \item \texttt{to\_bits()} 是位精确比较，浮点数 $0.0 \ne -0.0$；
  \item 命中已有索引则复用，否则分配新 \texttt{VertexId}。
\end{itemize}

\textbf{为何用位精确而非 $\epsilon$-tolerance？}
\begin{itemize}
  \item STL 文件中相同顶点的浮点表示通常\textbf{完全相同}（同一写入器输出）；
  \item 位精确比较 $O(1)$ 且无歧义；
  \item 浮点容差比较需 $O(\log V)$ 查找（KD-tree），且阈值敏感。
\end{itemize}

若 STL 文件含微小浮点误差导致同顶点不同表示，应先用
\texttt{weld\_vertices}（按距离阈值合并）做后处理。

\section{PLY 二进制小端读写}

PLY 二进制格式在文件大小与解析速度上显著优于 ASCII 变体，是工业应用的首选。
本模块支持完整的二进制小端格式（\texttt{format binary\_little\_endian 1.0}）。

\subsection{Header 解析}

PLY 文件由\textbf{文本 header} 与\textbf{二进制数据}两段组成，以 \texttt{end\_header\textbackslash n}
为分界。Header 描述了：
\begin{itemize}
  \item \texttt{format}：\texttt{ascii} / \texttt{binary\_little\_endian} / \texttt{binary\_big\_endian}；
  \item \texttt{element vertex <N>}：顶点元素数量；
  \item \texttt{property <type> <name>}：顶点属性（按声明顺序排列），支持
        \texttt{char} / \texttt{uchar} / \texttt{short} / \texttt{ushort} /
        \texttt{int} / \texttt{uint} / \texttt{float} / \texttt{double} 8 种类型；
  \item \texttt{element face <M>}：面元素数量；
  \item \texttt{property list <count\_type> <index\_type> <name>}：面索引 list。
\end{itemize}

\textbf{搜索算法}：用 \texttt{windows(11).position(|w| w == b"end\_header\textbackslash n")}
找到字节偏移，将 header 部分按文本逐行解析后，剩余字节按二进制布局读取。

\subsection{二进制数据布局}

顶点数据按 \texttt{vertex\_stride} 字节连续排列：
\[
  \text{vertex\_stride} = \sum_{i=1}^{k} \text{size}(\text{type}_i)
\]
其中 $k$ 为属性数量，$\text{size}(\cdot)$ 取 $\{1, 2, 4, 8\}$ 字节。
顶点 $j$ 的属性 $i$ 起始偏移：
\[
  \text{offset}(j, i) = j \cdot \text{vertex\_stride} + \sum_{l=1}^{i-1} \text{size}(\text{type}_l)
\]

面数据每条由 1 字节（\texttt{uchar}）顶点数 + $k \cdot \text{size}(\text{index\_type})$ 字节索引组成。

\subsection{ASCII / 二进制自动判别}

\texttt{parse\_ply\_bytes} 通过解析 header 中的 \texttt{format} 行决定后续路径：
\begin{itemize}
  \item \texttt{ascii} $\to$ 全文件转 UTF-8 后走 \texttt{parse\_ply\_ascii\_with\_header}；
  \item \texttt{binary\_little\_endian} $\to$ 从 \texttt{end\_header} 偏移走 \texttt{parse\_ply\_binary\_le}；
  \item \texttt{binary\_big\_endian} $\to$ 返回 \texttt{Unsupported} 错误。
\end{itemize}

\subsection{PlyType 类型分发}

\texttt{PlyType} 枚举对 8 种 PLY 标量类型提供：
\begin{itemize}
  \item \texttt{size()}：字节宽度；
  \item \texttt{read\_le\_as\_f64(bytes)}：小端读为 \texttt{f64}（顶点坐标）；
  \item \texttt{read\_le\_as\_u32(bytes)}：小端读为 \texttt{u32}（面索引）；
  \item \texttt{write\_le\_from\_f64(v)} / \texttt{write\_le\_from\_u32(v)}：小端写字节。
\end{itemize}

\subsection{format\_ply\_binary 输出约定}

输出固定布局（与 OpenMesh / MeshLab 兼容）：
\begin{itemize}
  \item \texttt{format binary\_little\_endian 1.0}；
  \item 顶点属性 \texttt{float x y z}（每顶点 12 字节）；
  \item 面索引 \texttt{property list uchar int vertex\_indices}（每面 $1 + 4k$ 字节）。
\end{itemize}

\subsection{PlyError}

\begin{center}
\renewcommand{\arraystretch}{1.18}
\begin{tabular}{@{}ll@{}}
  \toprule
  \textbf{变体} & \textbf{触发场景} \\
  \midrule
  \texttt{Io(std::io::Error)}     & 文件读写错误 \\
  \texttt{Parse\{line, msg\}}     & header 或数据解析失败（行号 + 描述） \\
  \texttt{Unsupported(String)}    & 不支持的 PLY 特性（如大端） \\
  \bottomrule
\end{tabular}
\end{center}

\section{OFF 读写}

OFF（Object File Format）是 Geomview 与许多学术工具使用的简化网格格式。

\subsection{格式约定}

\begin{lstlisting}
OFF
<vertex_count> <face_count> <edge_count>
x y z                # vertex 0
...
k v0 v1 ... vk-1     # face 0 (k = vertex count of face)
...
\end{lstlisting}

\begin{itemize}
  \item 第一行 \texttt{OFF} 关键字可选；部分文件首行为 \texttt{OFF 4 4 0}（关键字与计数同行）；
  \item \texttt{\#} 开头行视为注释；
  \item \texttt{edge\_count} 通常为 0，被忽略；
  \item 面顶点数 $k \ge 3$ 任意，自动按 fan 方式三角化（通过
        \texttt{build\_mesh\_from\_polygons}）；
  \item 索引越界（$v_i \ge \text{vertex\_count}$）返回 \texttt{Parse} 错误。
\end{itemize}

\subsection{首行解析流程}

\begin{center}
\begin{tikzpicture}[
  node distance=0.45cm,
  proc/.style={draw, rounded corners=2pt, minimum width=10cm, minimum height=0.6cm, align=center, font=\small},
  dec/.style={diamond, draw, aspect=2.4, fill=orange!12, inner sep=1pt, font=\small, align=center, minimum width=3.5cm},
  op/.style={-Stealth, thick},
  startend/.style={draw, rounded corners=10pt, minimum width=2.4cm, minimum height=0.6cm, font=\small, fill=gray!12}
]
  \node[startend] (s) {读取首行};
  \node[dec, below=of s] (d1) {首行以 \texttt{OFF} 起始？};
  \node[proc, below=of d1, fill=blue!8] (yes) {剥离 \texttt{OFF} 前缀，剩余作为计数行};
  \node[dec, below=of yes] (d2) {剩余非空？};
  \node[proc, right=1.6cm of d2, fill=green!12] (use) {用剩余作为计数行};
  \node[proc, below=of d2, fill=yellow!15] (next) {读取下一行作为计数行};
  \node[proc, below=of next, fill=blue!8] (no) {首行整体作为计数行};
  \node[startend, below=2cm of no, fill=gray!12] (e) {解析 V / F / E};
  \draw[op] (s) -- (d1);
  \draw[op] (d1) -- node[right, font=\scriptsize] {是} (yes);
  \draw[op] (yes) -- (d2);
  \draw[op] (d2) -- node[above, font=\scriptsize] {是} (use);
  \draw[op] (d2) -- node[right, font=\scriptsize] {否} (next);
  \draw[op] (d1) -- node[above, font=\scriptsize] {否} ++(3.6,0) |- (no.east);
  \draw[op] (use) |- (e);
  \draw[op] (next) -- (e);
  \draw[op] (no) -- (e);
\end{tikzpicture}
\end{center}

\subsection{OffError}

\begin{center}
\renewcommand{\arraystretch}{1.18}
\begin{tabular}{@{}ll@{}}
  \toprule
  \textbf{变体} & \textbf{触发场景} \\
  \midrule
  \texttt{Io(std::io::Error)}     & 文件读写错误 \\
  \texttt{Parse\{line, msg\}}     & 解析失败（行号 + 描述） \\
  \bottomrule
\end{tabular}
\end{center}

\section{glTF GLB 读写}

glTF（GL Transmission Format）是 Khronos 制定的现代 3D 资源交换标准。
本模块实现 GLB（glTF 二进制容器）的\textbf{最小子集}：

\begin{itemize}
  \item 仅支持单 mesh 单 primitive，mode = 4（TRIANGLES）；
  \item POSITION 属性必须为 \texttt{FLOAT}（componentType 5126）；
  \item 索引 accessor 可选 \texttt{UBYTE}（5121） / \texttt{USHORT}（5123） / \texttt{UINT}（5125），
        若无 indices 则按顺序 0..N；
  \item \textbf{不支持}材质、纹理、动画、skin、camera、node 层级等。
\end{itemize}

\subsection{GLB 文件结构}

\begin{center}
\begin{tikzpicture}[
  node distance=0.3cm,
  block/.style={draw, minimum width=12cm, minimum height=0.7cm, align=center, font=\small},
  op/.style={-Stealth, thick}
]
  \node[block, fill=gray!15] (h) {12 字节 GLB Header：magic=0x46546C67, version=2, total\_length};
  \node[block, below=of h, fill=blue!10] (jc) {JSON Chunk：4 字节长度 + 4 字节类型 0x4E4F534A + JSON 数据（4 字节对齐）};
  \node[block, below=of jc, fill=green!10] (bc) {BIN Chunk：4 字节长度 + 4 字节类型 0x004E4942 + 二进制数据（4 字节对齐）};
  \draw[op] (h) -- (jc);
  \draw[op] (jc) -- (bc);
\end{tikzpicture}
\end{center}

\subsection{JSON 与极简 JSON 解析器}

为避免引入 \texttt{serde\_json} 依赖，本模块实现一个\textbf{极简 JSON 解析器}
（\texttt{parse\_minimal\_json}），仅支持 GLB 中所需的对象/数组/原始值。
解析树 \texttt{JsonValue} 提供 \texttt{get(key)} / \texttt{as\_array()} / \texttt{as\_u32()}
等访问方法。

\subsection{GLB 解析流程}

\begin{center}
\resizebox{\textwidth}{!}{%
\begin{tikzpicture}[
  node distance=0.45cm,
  proc/.style={draw, rounded corners=2pt, minimum width=12cm, minimum height=0.65cm, align=center, font=\small},
  op/.style={-Stealth, thick},
  startend/.style={draw, rounded corners=10pt, minimum width=2.4cm, minimum height=0.65cm, font=\small, fill=gray!12}
]
  \node[startend] (s) {字节流};
  \node[proc, below=of s, fill=blue!8] (p1) {1. 校验 magic / version，读 total\_length};
  \node[proc, below=of p1, fill=blue!8] (p2) {2. 循环读 chunks：length + type + data，4 字节对齐};
  \node[proc, below=of p2, fill=blue!8] (p3) {3. JSON chunk $\to$ \texttt{parse\_minimal\_json} $\to$ \texttt{GltfDoc}};
  \node[proc, below=of p3, fill=blue!8] (p4) {4. 取 meshes[0].primitives[0]，查 POSITION 与 indices accessor};
  \node[proc, below=of p4, fill=blue!8] (p5) {5. 按 accessor 读取 BIN chunk 中的位置与索引};
  \node[proc, below=of p5, fill=blue!8] (p6) {6. mode=4 时按每 3 索引构成三角面};
  \node[proc, below=of p6, fill=yellow!15] (p7) {7. \texttt{build\_mesh\_from\_vertices\_and\_faces}};
  \node[startend, below=of p7] (e) {半边网格};
  \draw[op] (s) -- (p1);
  \draw[op] (p1) -- (p2);
  \draw[op] (p2) -- (p3);
  \draw[op] (p3) -- (p4);
  \draw[op] (p4) -- (p5);
  \draw[op] (p5) -- (p6);
  \draw[op] (p6) -- (p7);
  \draw[op] (p7) -- (e);
\end{tikzpicture}
}
\end{center}

\subsection{Accessor 读取}

\textbf{POSITION}（\texttt{read\_accessor\_f32x3}）：
\begin{itemize}
  \item 校验 \texttt{componentType == 5126}（FLOAT）；
  \item 计算起始偏移 $s = \text{bufferView.byteOffset} + \text{accessor.byteOffset}$；
  \item 每顶点 12 字节（3 × f32 LE），转 \texttt{f64} 推入 \texttt{Vec<[f64; 3]>}。
\end{itemize}

\textbf{indices}（\texttt{read\_accessor\_u32}）：
\begin{itemize}
  \item \texttt{componentType} $\in \{5121, 5123, 5125\}$（UBYTE / USHORT / UINT）；
  \item 按元素大小 1/2/4 字节读为 \texttt{u32}。
\end{itemize}

\subsection{format\_glb 输出约定}

输出固定结构（与 Babylon.js / Three.js 兼容）：
\begin{itemize}
  \item GLB header：\texttt{magic=0x46546C67, version=2}；
  \item JSON chunk：描述 asset / scene / nodes / meshes / buffers / bufferViews / accessors，
        JSON 文本用空格填充至 4 字节对齐；
  \item BIN chunk：先 $12V$ 字节 position（float32 × 3V），后 $4M$ 字节 indices（uint32 × M），
        末尾用 0 填充至 4 字节对齐；
  \item 仅三角面输出，非三角面跳过。
\end{itemize}

\subsection{GltfError}

\begin{center}
\renewcommand{\arraystretch}{1.18}
\begin{tabular}{@{}ll@{}}
  \toprule
  \textbf{变体} & \textbf{触发场景} \\
  \midrule
  \texttt{Io(std::io::Error)}     & 文件读写错误 \\
  \texttt{Parse(String)}          & GLB 结构或 JSON 解析失败 \\
  \texttt{Unsupported(String)}    & 不支持的 glTF 特性（如 mode $\ne 4$） \\
  \bottomrule
\end{tabular}
\end{center}

\section{统一入口：load\_mesh / save\_mesh / MeshError}

为简化上层调用，\texttt{io.rs} 提供\textbf{按扩展名自动分派}的统一接口：

\begin{center}
\renewcommand{\arraystretch}{1.18}
\begin{tabular}{@{}lll@{}}
  \toprule
  \textbf{扩展名} & \texttt{load\_mesh} 调用 & \texttt{save\_mesh} 调用 \\
  \midrule
  \texttt{.obj}              & \texttt{load\_obj}  & \texttt{save\_obj} \\
  \texttt{.ply}              & \texttt{load\_ply}（自动判别 ASCII / 二进制） & \texttt{save\_ply}（ASCII） \\
  \texttt{.stl}              & \texttt{load\_stl}（自动判别 ASCII / 二进制） & \texttt{save\_stl\_ascii} \\
  \texttt{.off}              & \texttt{load\_off}  & \texttt{save\_off} \\
  \texttt{.glb} / \texttt{.gltf} & \texttt{load\_glb} & \texttt{save\_glb} \\
  其他 & \multicolumn{2}{c}{\texttt{Err(MeshError::UnsupportedFormat(ext))}} \\
  \bottomrule
\end{tabular}
\end{center}

\subsection{MeshError 枚举}

\texttt{MeshError} 是 OBJ / PLY / STL / OFF / glTF 错误的\textbf{统一包装}，
通过 \texttt{From<ObjError>} / \texttt{From<PlyError>} / \texttt{From<StlError>}
/ \texttt{From<OffError>} / \texttt{From<GltfError>} 实现自动转换，
便于在 \texttt{?} 运算符中混用：

\begin{center}
\renewcommand{\arraystretch}{1.18}
\begin{tabular}{@{}ll@{}}
  \toprule
  \textbf{变体} & \textbf{来源} \\
  \midrule
  \texttt{Obj(ObjError)}                & \texttt{load\_obj} / \texttt{save\_obj} 错误 \\
  \texttt{Ply(PlyError)}                & \texttt{load\_ply} / \texttt{save\_ply*} 错误 \\
  \texttt{Stl(StlError)}                & \texttt{load\_stl} / \texttt{save\_stl\_*} 错误 \\
  \texttt{Off(OffError)}                & \texttt{load\_off} / \texttt{save\_off} 错误 \\
  \texttt{Gltf(GltfError)}              & \texttt{load\_glb} / \texttt{save\_glb} 错误 \\
  \texttt{UnsupportedFormat(String)}    & 未知扩展名 \\
  \bottomrule
\end{tabular}
\end{center}

\subsection{使用示例}

\begin{lstlisting}
// 自动按扩展名加载
let mesh = halfedge::load_mesh("model.obj")?;
let mesh = halfedge::load_mesh("model.ply")?;   // 自动判别 ASCII / 二进制
let mesh = halfedge::load_mesh("model.off")?;
let mesh = halfedge::load_mesh("model.glb")?;

// 自动按扩展名保存
halfedge::save_mesh(&mesh, "out.ply")?;          // ASCII
halfedge::save_mesh(&mesh, "out.glb")?;          // 二进制
halfedge::save_ply_binary(&mesh, "out_bin.ply")?; // 显式二进制 PLY
\end{lstlisting}

\section{测试覆盖}

\begin{longtable}{@{}ll@{}}
  \toprule
  \textbf{测试名} & \textbf{验证点} \\
  \midrule
  \endhead
  \multicolumn{2}{l}{\textit{构建器（10 个）}} \\
  \texttt{build\_mesh\_basic\_quad}                  & 4 顶点 2 面 10 半边 \\
  \texttt{build\_mesh\_passes\_full\_validation}     & 构建结果通过完整校验 \\
  \texttt{build\_mesh\_closed\_tetrahedron}          & 闭合四面体 12 半边 \\
  \texttt{build\_mesh\_empty\_inputs\_returns\_empty\_mesh} & 空顶点 + 空面 → 空网格，不 panic \\
  \texttt{build\_mesh\_vertices\_no\_faces}          & 仅顶点无面 → 顶点保留无面 \\
  \texttt{build\_polygons\_empty\_inputs\_returns\_empty\_mesh} & \texttt{build\_mesh\_from\_polygons} 空输入 \\
  \texttt{build\_mesh\_face\_index\_out\_of\_range\_panics} & 面索引越界 \texttt{panic!("越界")} \\
  \texttt{build\_polygons\_face\_index\_out\_of\_range\_panics} & 多边形入口同样越界 panic \\
  \texttt{build\_polygons\_skips\_degenerate\_face\_2\_verts} & 2 顶点退化面被静默跳过 \\
  \texttt{build\_polygons\_mixed\_degenerate\_and\_valid} & 退化与有效面混合时只保留有效面 \\
  \midrule
  \multicolumn{2}{l}{\textit{OBJ 往返与解析（12 个）}} \\
  \texttt{obj\_roundtrip\_quad}                      & 四边形 OBJ 往返一致 \\
  \texttt{obj\_roundtrip\_tetrahedron}               & 四面体 OBJ 往返一致 \\
  \texttt{obj\_parse\_skips\_comments\_and\_other\_lines} & 忽略 \texttt{\#} / \texttt{vt} / \texttt{vn} / \texttt{g} 等 \\
  \texttt{obj\_parse\_supports\_v\_vt\_vn\_format}   & 解析 \texttt{v/vt/vn} 格式 \\
  \texttt{obj\_parse\_negative\_indices}             & 负索引支持 \\
  \texttt{obj\_parse\_quadrilateral\_face\_succeeds} & 四边形面解析成功（支持 n-gon） \\
  \texttt{obj\_parse\_out\_of\_range\_index\_fails}  & 越界索引报 \texttt{IndexOutOfRange} \\
  \texttt{obj\_save\_load\_file\_roundtrip}          & 文件 IO 往返一致 \\
  \texttt{obj\_parse\_empty\_text}                   & 空文本 → 空网格，不报错 \\
  \texttt{obj\_parse\_only\_vertices\_no\_faces}     & 仅 \texttt{v} 行 → 顶点保留无面 \\
  \texttt{obj\_parse\_face\_zero\_index\_fails}      & 0 索引（OBJ 1-based）报错 \\
  \texttt{obj\_parse\_face\_index\_equal\_to\_vertex\_count\_fails} & 索引 = 顶点数报越界 \\
  \midrule
  \multicolumn{2}{l}{\textit{PLY 二进制（6 个）}} \\
  \texttt{ply\_binary\_roundtrip}                    & 二进制 PLY 字节往返一致 \\
  \texttt{ply\_binary\_file\_roundtrip}              & 二进制 PLY 文件 IO 往返 \\
  \texttt{ply\_ascii\_still\_works\_via\_parse\_ply\_bytes} & 旧 ASCII 文本经字节入口仍可解析 \\
  \texttt{ply\_binary\_detects\_bad\_header}         & 缺失 \texttt{end\_header} 报 \texttt{Parse} \\
  \texttt{ply\_parse\_empty\_bytes\_fails}           & 空字节流报错 \\
  \texttt{ply\_parse\_missing\_end\_header\_fails}   & 缺失 \texttt{end\_header} 报 \texttt{Parse} \\
  \midrule
  \multicolumn{2}{l}{\textit{STL 往返与解析（9 个）}} \\
  \texttt{stl\_ascii\_roundtrip}                     & STL ASCII 往返一致 \\
  \texttt{stl\_binary\_roundtrip}                    & STL 二进制往返一致 \\
  \texttt{stl\_ascii\_parses\_minimum\_solid}        & 最小 \texttt{solid} 块解析成功 \\
  \texttt{stl\_ascii\_rejects\_non\_triangular}      & 非三角面报 \texttt{NotTriangular} \\
  \texttt{stl\_binary\_detects\_size\_mismatch}      & 二进制长度不匹配报 \texttt{BadBinarySize} \\
  \texttt{stl\_file\_roundtrip\_ascii}               & STL ASCII 文件 IO 往返 \\
  \texttt{stl\_file\_roundtrip\_binary}              & STL 二进制文件 IO 往返 \\
  \texttt{stl\_ascii\_empty\_text}                   & 空文本 → 空网格，不报错 \\
  \texttt{stl\_ascii\_only\_solid\_endsolid}         & 仅 \texttt{solid/endsolid} 块 → 空网格 \\
  \midrule
  \multicolumn{2}{l}{\textit{OFF 往返与解析（7 个）}} \\
  \texttt{off\_roundtrip\_tetrahedron}               & 四面体 OFF 往返一致 \\
  \texttt{off\_file\_roundtrip}                      & OFF 文件 IO 往返 \\
  \texttt{off\_parse\_counts\_inline\_with\_keyword} & 首行 \texttt{OFF 4 4 0} 行内计数支持 \\
  \texttt{off\_parse\_quadrilateral\_face\_succeeds} & 四边形面解析成功 \\
  \texttt{off\_parse\_out\_of\_range\_index\_fails}  & 索引越界报 \texttt{Parse} \\
  \texttt{off\_parse\_empty\_text\_fails}            & 空文本报错（OFF 需首行关键字） \\
  \texttt{off\_parse\_only\_vertices\_no\_faces}     & 仅顶点无面 → 顶点保留无面 \\
  \midrule
  \multicolumn{2}{l}{\textit{glTF GLB 往返与解析（4 个）}} \\
  \texttt{glb\_roundtrip\_tetrahedron}               & 四面体 GLB 字节往返一致 \\
  \texttt{glb\_file\_roundtrip}                      & GLB 文件 IO 往返 \\
  \texttt{glb\_detects\_bad\_magic}                  & 错误 magic 报 \texttt{Parse} \\
  \texttt{glb\_icosphere\_roundtrip}                 & icosphere 经 GLB 往返一致 \\
  \midrule
  \multicolumn{2}{l}{\textit{统一入口（6 个）}} \\
  \texttt{auto\_detect\_obj}                         & \texttt{load\_mesh(".obj")} 走 OBJ 路径 \\
  \texttt{auto\_detect\_ply}                         & \texttt{load\_mesh(".ply")} 走 PLY 路径 \\
  \texttt{auto\_detect\_stl}                         & \texttt{load\_mesh(".stl")} 走 STL 路径 \\
  \texttt{auto\_detect\_off}                         & \texttt{load\_mesh(".off")} 走 OFF 路径 \\
  \texttt{auto\_detect\_glb}                         & \texttt{load\_mesh(".glb")} 走 GLB 路径 \\
  \texttt{auto\_detect\_unsupported}                 & 未知扩展名返回 \texttt{UnsupportedFormat} \\
  \midrule
  \multicolumn{2}{l}{\textit{损坏输入与退化面（11 个，散布于上述分组）}} \\
  \multicolumn{2}{l}{上述 \texttt{*\_empty\_*} / \texttt{*\_out\_of\_range\_*} / \texttt{*\_degenerate\_*} / \texttt{*\_only\_*} 系列} \\
  \multicolumn{2}{l}{覆盖空输入、越界索引、退化面、仅顶点等边界场景，} \\
  \multicolumn{2}{l}{确保所有解析入口在异常输入下既不 \texttt{unwrap} panic、也不静默吞错。} \\
  \bottomrule
\end{longtable}

\subsection{历史 panic 回归}

\texttt{tests/panic\_safety.rs} 集成测试中针对 \texttt{io} 模块的回归用例：

\begin{sloppypar}
\begin{itemize}
  \item \texttt{regression\_build\_mesh\_index\_out\_of\_range\_panics\_with\_clear\_message}：保证越界 \texttt{panic!} 消息含中文 \texttt{越界}；
  \item \texttt{regression\_parse\_obj\_empty\_input\_no\_panic}：空 OBJ 文本不 panic；
  \item \texttt{regression\_parse\_obj\_only\_vertices\_no\_faces\_no\_panic}：仅顶点 OBJ 不 panic；
  \item \texttt{regression\_parse\_obj\_malformed\_face\_skipped\_no\_panic}：损坏面行不 panic。
\end{itemize}
\end{sloppypar}

全模块共 \textbf{54 个}单元测试，全库共 \textbf{578 个}单元测试 + \textbf{24 个}集成测试（\texttt{tests/mesh\_workflow.rs} 与 \texttt{tests/panic\_safety.rs}）。

\end{document}
